\section{事故一：Retry、Cron 与 Rebuild 的正反馈}

课堂第一段事故不是“一个数据库挂了”，而是多个自动机制同时响应：timeout 触发 retry，repair job 和 cron 增加负载，queue 拉长，cache miss 与 rebuild 又进一步压垮依赖。Incident response 的首要动作不是添加更多修复工作，而是识别正反馈并停止 producer。

\subsection{从 initial fault 到 cascade}

本节把课堂事故改写成 dependency feedback graph，而不是按时间复述操作。读图时要从 initial fault 追踪 automatic reaction、load amplification 和 mitigation：关键判断不是哪个组件先报错，而是哪些 retry、repair job 和 rebuild 在持续增加 saturated dependency 的输入，以及哪个控制点能最快降低系统增益。

\lecturefigure{07-incident-cascade.png}{自动 retry 和 repair job 可能把局部故障放大成全局事故。}{本地字幕 19:10--25:30；概念重绘。}

\begin{knowledgebox}{读图：恢复顺序与平时优化顺序相反}
平时追求吞吐和自动化；事故中先保护 invariant：停止非关键 producer、限制 retry、shed load、冻结 deploy，再恢复最小 read/write path。只有 dependency 有 headroom 后，才逐步放开 backlog 和 rebuild。
\end{knowledgebox}

\begin{importantbox}{Retry amplification}
若原始请求率为 $\lambda$，每次失败平均触发 $r$ 次重试，失败概率为 $p$，附加负载可近似为
\[
\lambda_{effective}\approx \lambda(1+pr+(pr)^2+\cdots).
\]
当 $pr$ 接近 1 时，重试链会非常敏感。实际系统用 bounded retry、exponential backoff、jitter、circuit breaker 和 global retry budget 控制放大。
\end{importantbox}

\teachervoice{课程中的事故叙述有很强的“一个人坐在那里救系统”的戏剧感。真正应该学习的是 command structure：谁能停止 producer、谁判断数据 invariant、谁负责用户沟通，以及哪些自动化在 degraded mode 必须关闭。}

\subsection{本章小结}

Incident cascade 的根因常是 feedback loop，而非单点组件。恢复需要先降低系统增益，再修复状态；postmortem 应把 retry、cron、queue 与 migration 纳入同一 dependency graph。

